\begin{table}[htbp]
\centering
\small
\setlength{\tabcolsep}{3pt}
\renewcommand{\arraystretch}{1.08}
\begin{tabularx}{\textwidth}{C{0.8cm} Y Y Y Y}
\hline
\textbf{Step} &
\textbf{Input} &
\textbf{Processing Stage} &
\textbf{Output} &
\textbf{Next Module} \\
\hline
1 &
Workpiece geometry, bend-side selection, \ac{tcp}, gripper, and safety parameters &
Configuration and operator-input initialisation &
Structured task and configuration data &
Safety validation \\
\hline
2 &
Configuration data, fixture states, and operator start decision &
Pre-execution validation and runtime readiness checks &
Permission to continue or controlled stop &
Pose acquisition and zone classification \\
\hline
3 &
Tracked sheet pose and workspace zone &
Orientation admissibility and pickup-side validation &
Selected pickup orientation or conservative fallback &
Grip planning \\
\hline
4 &
Selected orientation, sheet geometry, and gripper parameters &
Grip candidate generation, filtering, and force validation &
Feasible grip configuration and pickup poses &
Regrip and sequence planning \\
\hline
5 &
Current workpiece pose, selected grip, and target bending pose &
Direct reachability, regrip, and opposite-side handling logic &
Validated task sequence with required handling actions &
Trajectory generation \\
\hline
6 &
Ordered task poses and handling actions &
Interpolation, kinematic validation, path reduction, and command formatting &
Executable robot motion and gripper commands &
AnyRobot runtime communication \\
\hline
\end{tabularx}
\caption{End-to-end data flow from workpiece and process input to executable robot commands.}
\label{tab:ch5-end-to-end-data-flow}
\end{table}
